\begin{resumo}
\begin{SingleSpacing}
A conservação contínua da infraestrutura urbana constitui um desafio estrutural para a administração pública de municípios brasileiros de médio e grande porte, nos quais falhas cotidianas --- como buracos na pavimentação asfáltica, lâmpadas de iluminação pública inoperantes, obstruções ou vazamentos em redes de drenagem pluvial, semáforos defeituosos e pontos de descarte irregular de resíduos sólidos --- comprometem a segurança viária, a salubridade ambiental e o bem-estar coletivo. Na maioria das prefeituras, o fluxo informacional entre o cidadão que vivencia o problema no território e as equipes operacionais encarregadas do reparo caracteriza-se pela fragmentação, gerando chamados duplicados, imprecisão na localização geográfica das ocorrências, ausência de evidências visuais padronizadas e falta de controle sobre o despacho de frota, a alocação de mão de obra e o consumo de insumos de almoxarifado. Diante desse cenário, este relatório técnico-acadêmico, referente à Etapa 1 do projeto da disciplina de Banco de Dados, tem como objetivo descrever detalhadamente o problema da gestão colaborativa de zeladoria e manutenção de infraestrutura urbana, fundamentar seu enquadramento nos Objetivos de Desenvolvimento Sustentável (ODS) da Agenda 2030 da Organização das Nações Unidas (ONU) e delimitar o escopo funcional de um sistema de banco de dados relacional voltado à sua resolução. A análise apoia-se em dados oficiais do Instituto Brasileiro de Geografia e Estatística (IBGE), do Instituto de Pesquisa Econômica Aplicada (IPEA), do Sistema Nacional de Informações sobre Saneamento (SNIS) e da legislação nacional de gestão urbana e defesa do usuário do serviço público. Demonstra-se que a proposta se alinha diretamente ao ODS 11 (Cidades e Comunidades Sustentáveis), com ênfase nas Metas 11.3, 11.6 e 11.7, e de modo complementar ao ODS 9 (Indústria, Inovação e Infraestrutura, Meta 9.1) e ao ODS 16 (Paz, Justiça e Instituições Eficazes, Metas 16.6 e 16.7). Por fim, estabelece-se o escopo de alto nível do banco de dados composto por doze entidades de negócio organizadas em três subdomínios integrados --- demandas cidadãs geolocalizadas por CEP, estrutura operacional das secretarias municipais e execução de ordens de serviço articulada à cadeia de suprimentos ---, preparando a base conceitual para a modelagem entidade-relacionamento e relacional desenvolvida na Etapa 2.

\vspace{\onelineskip}
\noindent
\textbf{Palavras-chave}: Infraestrutura urbana. Zeladoria municipal. Banco de dados relacional. Objetivos de Desenvolvimento Sustentável. ODS 11.
\end{SingleSpacing}
\end{resumo}
